% from ega4-11-fr.tex
\oldpage[IV-3]{127}

En effet, par récurrence sur $k$, l'hypothèse que les
$\mathfrak J^kM/\mathfrak J^{k+1}M$ sont des $A$-modules plats pour $k\geq0$
entraîne d'abord, par ($\mathbf 0_{\mathrm I}$, 6.1.2), que
$M/\mathfrak J^{k+1}M$ est un $A$-module plat~; en outre
($\mathbf 0_{\mathrm I}$, 6.1.2), la suite
\[
0\longrightarrow(\mathfrak J^{k+1}M)\otimes_AA'
\longrightarrow M\otimes_AA'
\longrightarrow(M/\mathfrak J^{k+1}M)\otimes_AA'\longrightarrow0
\]
est alors exacte, autrement dit $(\mathfrak J^{k+1}M)\otimes_AA'$ s'identifie
à son image canonique $\mathfrak J'^{,k+1}M'$ dans $M'=M\otimes_AA'$.
D'autre part, toujours en vertu de ($\mathbf 0_{\mathrm I}$, 6.1.2), la suite
\[
0\longrightarrow(\mathfrak J^{k+1}M)\otimes_AA'
\longrightarrow(\mathfrak J^kM)\otimes_AA'
\longrightarrow(\mathfrak J^kM/\mathfrak J^{k+1}M)\otimes_AA'
\longrightarrow0
\]
est exacte, ce qui prouve le lemme.

Pour démontrer la nécessité des conditions de (11.2.9), nous procéderons en
plusieurs étapes.

\medskip
\noindent\textbf{I)} \emph{Réduction au cas où les $A_\lambda$ sont
noethériens.} --- On procède comme dans la réduction I) de (11.2.6.1), dont
nous gardons les notations~; il faut simplement commencer par remplacer
$A'_0$ par une sous-$A'_0$-algèbre de type fini $A''_0$ de $A_0$ telle que,
si l'on pose $B''_0=B'_0\otimes_{A'_0}A''_0$, il y ait dans $B''_0$ un idéal
de type fini $\mathfrak J''_0$ tel que $\mathfrak J''_0B_0=\mathfrak J_0$.
On considère pour cela les sous-$A'_0$-algèbres $A'_\alpha$ de type fini de
$A_0$, qui forment une famille filtrante, et l'on a
$B_0=\varinjlim B'_\alpha$, où
$B'_\alpha=B'_0\otimes_{A'_0}A'_\alpha$~; il y a donc un indice $\beta$ tel
qu'un système fini de générateurs de $\mathfrak J_0$ soit formé d'images dans
$B_0$ d'éléments de $B'_\beta$~; on prendra alors $A''_0=A'_\beta$,
$B''_0=B'_\beta$ et pour $\mathfrak J''_0$ l'idéal engendré par ces éléments.
On peut donc supposer que $A_0$ (donc aussi $B_0$) est noethérien. On définit
ensuite comme \emph{loc. cit.} l'ensemble filtrant $H$ et les $A_\xi$,
$B_\xi$, $M_\xi$ pour $\xi\in H$~; on posera aussi
$\mathfrak J_\xi=\mathfrak J_0B_\xi$ pour tout $\xi\in H$. Supposons alors
que l'on ait démontré qu'il existe un $\xi=(\lambda,C_\lambda)$ tel que
$\operatorname{gr}_{\mathfrak J_\xi}^*(M_\xi)$ soit un $C_\lambda$-module
plat~; comme $M_\lambda=M_\xi\otimes_{C_\lambda}A_\lambda$, il résulte de
(11.2.9.2) que $\operatorname{gr}_{\mathfrak J_\lambda}^*(M_\lambda)$ est un
$A_\lambda$-module plat.

\medskip
\noindent\textbf{II)} \emph{Lemmes préliminaires.}

\begin{lemma}[11.2.9.3]
\label{IV.11.2.9.3-fr}
Soient $A$ un anneau,
$B=\bigoplus_{i\geq0}B_i$ une $A$-algèbre graduée à degrés positifs,
$M=\bigoplus_{i\in\mathbf Z}M_i$ un $B$-module gradué. On suppose que $B_0$
est un anneau local, et que chacun des $M_i$ est un $B_0$-module de type fini.
Pour que $M$ soit un $B$-module de type fini, il faut et il suffit que, si
$\mathfrak q$ est l'image réciproque dans $A$ de l'idéal maximal $\mathfrak m$
de $B_0$, $M\otimes_Ak(\mathfrak q)$ soit un
$(B\otimes_Ak(\mathfrak q))$-module de type fini.
\end{lemma}

La condition étant évidemment nécessaire, prouvons qu'elle est suffisante.
Comme $B_0$ (et par suite $B$) est une $A_{\mathfrak q}$-algèbre, on peut
remplacer $A$ par $A_{\mathfrak q}$, autrement dit supposer que $A$ est aussi
un anneau local dont $\mathfrak q$ est l'idéal maximal. Par hypothèse, il existe
un entier $N$ tel que l'homomorphisme canonique
\[
\left(\bigoplus_{-N\leq i\leq N}(M_i\otimes_{B_0}B)\right)
\otimes_Ak(\mathfrak q)\longrightarrow M\otimes_Ak(\mathfrak q)
\]

